\section{总结}

本文建立了由太阳位置、双向量定日镜跟踪、光线级阴影遮挡与截断判定组成的统一光学模型，并据此完成给定镜场评价、统一规格镜场优化和逐镜连续规格优化。

对于问题一，给定镜场在题目60个太阳状态下的年平均光学效率为0.5723，年平均输出热功率为34.9868 MW，单位镜面面积年平均输出热功率为0.5569 kW/m$^2$。对于问题二，最密三角晶格与分区参数化搜索得到满足60 MW约束的统一规格方案，其年平均输出热功率为60.0257 MW，单位面积功率为0.4472 kW/m$^2$。对于问题三，从空场重新生成布局并连续调整逐镜宽度、高度和安装高度后，最终2939面定日镜的总面积为132952.1202 m$^2$，年平均输出热功率为60.1119 MW，单位面积功率达到0.4521 kW/m$^2$，较问题二提高1.1127\%。

综合三问可见，在相同光学评价口径下，参数化密排能够控制大规模镜位搜索的计算量，逐镜连续规格则能进一步削减低贡献镜面面积。最终结果均通过分辨率和多随机种子复核，但其有效范围仍受离散太阳状态、理想化太阳盘以及未计入工程误差等假设限制，因而应将其理解为本文模型与搜索预算下的可行设计结果。
